\chapter{Witten-Donaldson Theory And The Seiberg-Witten Comparison}
\label{ch:tqft-witten-donaldson-seiberg-witten-comparison}

\ConventionsEuclidean

This chapter begins the four-dimensional cohomological gauge theory that
connects Donaldson invariants to the low-energy Seiberg-Witten description.
The comparison is treated as an RG statement about a twisted
\(\mathcal N=2\) gauge theory and its Abelian Coulomb-branch effective theory.

\section{Twisted Yang-Mills Datum}

Let \(X\) be a compact oriented smooth four-manifold and let \(P\to X\) be a
principal \(SU(2)\) bundle.  The Donaldson twist of pure
\(\mathcal N=2\) Yang--Mills theory uses the scalar supercharge \(Q\) of
Chapter~\ref{ch:tqft-twists-of-supersymmetric-theories}.  The bosonic gauge
field is a connection \(A\) on \(P\), and the localization equation in the
ultraviolet weak-coupling chart is
\[
  F_A^+=0 .
\]
The instanton number is
\[
  k=-\frac{1}{8\pi^2}\int_X\operatorname{tr}(F_A\wedge F_A)
\]
in the trace convention used for the gauge theory.

\section{Instanton Moduli Space}

The anti-self-dual moduli space is
\[
  \mathcal M_k(P)
  =
  \{A\mid F_A^+=0\}/\mathcal G_P ,
\]
where \(\mathcal G_P\) is the gauge transformation group.  At an irreducible
smooth point, its tangent and obstruction spaces are the cohomology of
\[
  0\longrightarrow \Omega^0(X,\operatorname{ad}P)
  \overset{d_A}{\longrightarrow}
  \Omega^1(X,\operatorname{ad}P)
  \overset{d_A^+}{\longrightarrow}
  \Omega^{2,+}(X,\operatorname{ad}P)
  \longrightarrow0 .
\]
For \(SU(2)\), the virtual dimension is
\[
  \dim_{\mathrm{vir}}\mathcal M_k
  =
  8k-3(1-b_1(X)+b_2^+(X)).
\]
Compactification, reducible connections, orientations, and chamber structure
are part of the definition of the invariant.

\section{Donaldson Observables}

Let
\[
  \mathcal O_0=\operatorname{tr}\phi^2
\]
be the scalar observable in the twisted vector multiplet.  Descent gives
operators \(\mathcal O_j\) obeying
\[
  d\mathcal O_j=\{Q,\mathcal O_{j+1}\}_{\mathrm{gr}} .
\]
For a homology cycle \(\gamma_j\in H_j(X)\), define
\[
  I(\gamma_j)=\int_{\gamma_j}\mathcal O_j .
\]
The correlation functions of these insertions localize, in the ultraviolet
Donaldson chamber, to intersection pairings on the compactified instanton
moduli space when the degree of the insertion matches
\(\dim_{\mathrm{vir}}\mathcal M_k\).

\section{Coulomb-Branch RG Picture}

The same twisted theory has a low-energy description on the Coulomb branch.
For gauge algebra \(\mathfrak{su}(2)\), the coordinate \(u\) and periods
\((a_D,a)\) are those of
Chapter~\ref{ch:susy-four-dimensional-n2-seiberg-witten}.  Away from
singular values of \(u\), the low-energy theory is an Abelian twisted
\(\mathcal N=2\) gauge theory with coupling \(\tau(u)\).  At singular values,
monopole or dyon hypermultiplets become light and must be included.  The
topologically twisted path integral is therefore decomposed into a Coulomb
branch integral and local contributions from the singular fibers.

\section{Seiberg-Witten Monopole Equations}

Choose a \(\operatorname{Spin}^c\) structure on \(X\) with spinor bundles
\(S^\pm\) and determinant line \(L\).  Let \(A\) be a connection on \(L\) and
let \(q\in\Gamma(S^+\otimes L)\).  The Seiberg-Witten equations are
\[
  \not D_A q=0,
  \qquad
  F_A^+=\sigma(q),
\]
where \(\sigma(q)\in\Omega^{2,+}(X)\) is the quadratic map obtained from
Clifford multiplication.  These equations are the localization equations of
the twisted Abelian monopole effective theory near a singular Coulomb-branch
point.

\section{Comparison Mechanism}

The Donaldson and Seiberg-Witten descriptions are compared by the following
chain of objects:
\[
\begin{aligned}
  &\text{twisted non-Abelian UV theory}\\
  &\qquad\longrightarrow \text{Coulomb-branch Wilsonian theory}\\
  &\qquad\longrightarrow \text{twisted Abelian monopole theory}.
\end{aligned}
\]
The first arrow is an RG statement.  The second arrow is the local
Seiberg-Witten effective description near a singular fiber.  Contact terms,
metric dependence, \(u\)-plane contributions, wall crossing, and the map from
Donaldson insertions to Abelian observables must be included in the same
coordinate convention.

\section{Gap Ledger}

The finite-dimensional geometry of the instanton and monopole moduli spaces
can be developed with rigorous differential geometry once compactifications
and orientations are fixed.  The QFT comparison additionally requires a
construction of the twisted four-dimensional \(\mathcal N=2\) gauge theory,
a derivation of its Coulomb-branch Wilsonian action, control of singular
fiber contributions, and a proof that the cohomological correlation functions
are invariant along the RG comparison path.

\begin{openproblem}[Donaldson and Seiberg-Witten comparison as QFT]
\label{op:donaldson-seiberg-witten-qft-comparison}
Turn the physical RG comparison between Donaldson and Seiberg-Witten
invariants into a theorem-level statement of four-dimensional QFT, including
the ultraviolet twisted gauge theory, low-energy Abelian monopole theory,
observable map, contact terms, wall crossing, and compactification data.
\end{openproblem}
